%%%%%%%%%%%%%%%%%%%%%%%%%%%%%%%%%%%%%%%%%
% "ModernCV" CV and Cover Letter
% LaTeX Template
% Version 1.1 (9/12/12)
%
% This template has been downloaded from:
% http://www.LaTeXTemplates.com
%
% Original author:
% Xavier Danaux (xdanaux@gmail.com)
%
% License:
% CC BY-NC-SA 3.0 (http://creativecommons.org/licenses/by-nc-sa/3.0/)
%
% Important note:
% This template requires the moderncv.cls and .sty files to be in the same 
% directory as this .tex file. These files provide the resume style and themes 
% used for structuring the document.
%
%%%%%%%%%%%%%%%%%%%%%%%%%%%%%%%%%%%%%%%%%

%----------------------------------------------------------------------------------------
%	PACKAGES AND OTHER DOCUMENT CONFIGURATIONS
%----------------------------------------------------------------------------------------

\documentclass[10pt,a4paper]{moderncv} % Font sizes: 10, 11, or 12; paper sizes: a4paper, letterpaper, a5paper, legalpaper, executivepaper or landscape; font families: sans or roman

\moderncvstyle{classic} % CV theme - options include: 'casual' (default), 'classic', 'oldstyle' and 'banking'
\moderncvcolor{grey} % CV color - options include: 'blue' (default), 'orange', 'green', 'red', 'purple', 'grey' and 'black'
\usepackage{lipsum}
\usepackage{amsmath}
\usepackage{url}
\usepackage[scaled]{beramono}
\renewcommand*\familydefault{\ttdefault} %% Only if the base font of the document is to be typewriter style
\usepackage[T1]{fontenc}

\usepackage[scale=0.87]{geometry} % Reduce document margins
\linespread{1.1}
\setlength{\hintscolumnwidth}{2.16cm} % Uncomment to change the width of the dates column
%\setlength{\makecvtitlenamewidth}{10cm} % For the 'classic' style, uncomment to adjust the width of the space allocated to your name

%----------------------------------------------------------------------------------------
%	NAME AND CONTACT INFORMATION SECTION
%----------------------------------------------------------------------------------------

\firstname{K{\i}van\c{c}} % Your first name
\familyname{Tatar} % Your last name

% All information in this block is optional, comment out any lines you don't need
\title{Assistant Professor in Interactive AI}
\address{Chalmers University of Technology \\Data Science and AI Division \\ Computer Science Department \\ Rännvägen 6B,\\}{  
412 58 Gothenburg, Sweden}

\email{tatar@chalmers.se}
\homepage{aicomparts.com}{Research Group} 
\webpage{kivanctatar.com}{Personal Website} 

%----------------------------------------------------------------------------------------
\begin{document}
	%\ttfamily
	%\asciifamily
	\makecvtitle % Print the CV title
\begin{center} \LARGE  Description of Time Distribution\end{center}
\hline

\cvitem{}{My contract is on 80\% research and 20\% teaching. My research time is funded by WASP-HS.}
\cvitem{}{\textbf{I have written, prepared, and submitted my Docent and Associate Professor application} to the CSE department, which is now approved by the department to be sent to the appointment committee at Chalmers central. Writing my docent application and the pedagogical portfolio took significant amount of my time in 2025.}

\begin{center} \LARGE  Research Experience and Merits\end{center}
\hline

    \section{Research Grants}

    \cvitem{Summary}{I have acquired two significant grants as co-PI, totalling to a 48 million SEK. It is also important to note that I am the only WASP-HS assistant prof. who has acquired a WASP-HS cluster grant as a co-PI or PI.}
    
    {\begin{center}{\textcolor{magenta}{Submitted Grants Applications Still Awaiting Decisions}}\end{center}}

    \cvitem{ERC-1}{\href{https://erc.europa.eu/apply-grant/starting-grant}{European Research Council, Starting Grant}. \textit{Multimodal Generative Modelling for Interpretable Movement-to-Audio Synthesis: MODA}. Kıvanç Tatar (PI, Chalmers University of Technology and University of Gothenburg). 1.5 million Euro $\approx$16.4 million SEK.}

    
    \cvitem{}{\begin{center}\textcolor{magenta}{\textcolor{magenta}{Successful Grants}}\end{center}}
    

    \cvitem{}{\begin{center}\textcolor{magenta}{\textcolor{magenta}{Rejected Applications}}\end{center}}


\section{Academic Publications}

\cvitem{Summary}{\textbf{I have published two journal papers, and an additional journal paper is accepted and in press.} Additionally, three journal papers are submitted. It is exceptional that I have submitted 5 journal papers in 2025.}
    
	\cvitem{}{\begin{center} \textcolor{magenta}{\textcolor{magenta}{Journal Papers}}\end{center}}

    FDG PAPER
    AI and Society paper
    Hugh's resubmissions?

    \cvitem{1}{Zappi, V.,\textcolor{magenta}{Tatar, K}. (2025). \textit{Neural audio instruments: Epistemological and phenomenological perspectives on musical embodiment of deep learning}. \textbf{Frontiers in Computer Science}, 7. \url{https://doi.org/10.3389/fcomp.2025.1575168}}

    \cvitem{2}{Cotton, K., Kaila, A.K.,Jääskeläinen, P., Holzapel, A., \textcolor{magenta}{Tatar, K}. (2025).\textit{Imploding between the Facts and Concerns of Musical-AI: A Multidisciplinary Method for Analysing Human-AI Musical Interaction.} \textbf{Humanities and Social Sciences Communications}, 12(1), 1–20, \textbf{Springer Nature Press}. \url{https://doi.org/10.1057/s41599-025-04533-4}}

    \cvitem{3}{Madaghiele V., Lund L., Holzer D., Kelkar T., \textcolor{magenta}{Tatar, K.}, and Holzapfel A. \textit{Expanding the machine: notating generative synthesis with a state-based representation and an interactive timbre space}. \textbf{Accepted for publications and in press} at Organised Sound, Cambridge Press.}

    \cvitem{4}{Cotton, K., de Vries K., Holzapfel A., Berglund K., \textcolor{magenta}{Tatar, K.}. \textit{Revising MIR Research Practices for Singing Data Collection}. Submitted to AI and Society, Springer Nature.}

    \cvitem{5}{Liu X., Cotton K., \textcolor{magenta}{Tatar, K.} \textit{Revisiting a Decade of Computational Game Creativity}. Submitted to International Journal on Human Computer Interaction, Taylor and Francis. Passed the editor and now in review.}

    \cvitem{6}{*Fröhlich M. M., Cotton K., \textcolor{magenta}{Tatar, K.} \textit{Artificial Intelligence: Perspectives on a Caesura in Art, its Creation and Consumption}. Submitted to Leonardo, MIT Press.}    

	\cvitem{}{\small All papers are double-blind peer-reviewed unless otherwise is indicated.}
	\cvitem{}{\small *These papers are single-blind peer-reviewed. The reviewers had access to the list of authors.}

\begin{center} \LARGE  Teaching Experience and Merits\end{center}
\hline
    
\section{Supervision Experience}

\cvitem{Summary}{Kelsey Cotton has passed licentiate defense.}

\cvitem{}{\begin{center}\textcolor{magenta}{\textcolor{magenta}{Post-doctoral Fellows}}\end{center}}

\cvitem{Ongoing}{\href{https://scholar.google.com/citations?user=ni5MM4oAAAAJ}{\textbf{Xuechen (Hugh) Liu}}. \textit{Postdoc in Interactive AI for interdisciplinary Artistic Practices}. 2023-01::2026-01. \textcolor{magenta}{Supervisor: K{\i}van\c{c} Tatar}. }
\cvitem{}{\begin{center}\textcolor{magenta}{\textcolor{magenta}{Ph.D. students}}\end{center}}

\cvitem{Ongoing}{\href{https://kelseycotton.com/}{\textbf{Kelsey Cotton}}. \textit{Interactive Music Systems using Artificial Intelligence}. Expected graduation 2027. \textcolor{magenta}{Main Supervisor: K{\i}van\c{c} Tatar}, Co-supervisor: Devdatt Dubhashi, Examiner: Graham Kemp. Funded by WASP-HS-1. Licentiate defense is scheduled to 2024-11-26 with the external examiner Prof. Alexander Refsum Jensenius from University of Oslo, Norway.}

\cvitem{}{\textbf{New PhD}. \textbf{Changheon Han}. \textit{Generative AI in Culture and Creative Domains}. PhD hiring is ongoing. We are now interviewing shortlisted candidates. \textcolor{magenta}{Main Supervisor: K{\i}van\c{c} Tatar}, Co-supervisor: Ashkan Panahi, Examiner: Dag Wedelin.}

\cvitem{}{\begin{center}\textcolor{magenta}{\textcolor{magenta}{Master students}}\end{center}}

\cvitem{Summary}{}



\section{Teaching Experience}

\cvitem{2026 Spring}{\textbf{Examiner.} Advanced course. TRA 385 - Machine Learning and AI through artistic innovation, Chalmers University of Technology, Gothenburg, Sweden. As a part of \href{https://www.chalmers.se/en/education/your-studies/course-selection-and-registration/select-courses/choose-a-tracks-course/}{TRACKS} initiative, conducted using the \href{https://www.chalmers.se/en/education/study-at-chalmers/why-chalmers/fuse-chalmers-makerspace/}{FUSE} maker spaces of Chalmers. }

\cvitem{New Course?}{Masters in Acoustics program approaches me to initiate a new course in Musical AI and Music Information Retrieval. I have been in communication with DSAI teaching management (Birgit Grohe and Ashkan Panahi), who have been positive to this course. Ideally, this course will be an elective for both Masters in DSAI and Masters in Acoustics programs, thus reaching to a minimum of two masters programs. The course can run in 2028, so there is still time to plan.}

\section{Personal Development}

\cvitem{}{Intro to Quantum Computing: in-progress}

\begin{center} \LARGE Academic Citizenship\end{center}
\hline

\section{Activities at Chalmers}

\cvitem{2025-10-22}{\textit{TRA 385 - Machine Learning and AI through Artistic Innovation}. I presented the pedagogical approach and the curriculm design of TRA 385 on TRACKs Teachers' Day in 2025.}

\section{Reviewing Activities for Journals and Conferences}

FDG TWO LONG PAPERS
Journal paper reviews?

\section{Other}

\cvitem{}{I have joined four ERC information seminars and workshops covering the ERC Starting grant.}


\begin{center} \LARGE Research Utilization\end{center}
\hline

\section{Artworks exhibited in Public Venues}


\section{Press Coverage}


\end{document}